
\begin{table}[H]
\centering
\caption{Convergence in Beliefs}
\label{tab:4}
\begin{tabular}{lcc}
\toprule
\textbf{} & \multicolumn{2}{c}{\textbf{Inflation Expectations}}\\ 
\cmidrule(lr){2-3}
& (1) & (2) \\ 
\midrule
\multicolumn{3}{l}{\textbf{Relative to original COPOM statement (Treatment 1)}} \\
T2 & -0.714*** & -0.754*** \\
& (0.258) & (0.261) \\
\midrule
\multicolumn{3}{l}{\textbf{Relative to original COPOM statement (Treatment 1)}} \\
T3 & -0.264 & -0.414 \\
& (0.280) & (0.302) \\
\midrule
\multicolumn{3}{l}{\textbf{Relative to original G1 article (Treatment 2)}} \\
T4 & -0.197 & -0.182 \\
& (0.225) & (0.234) \\
\midrule
\multicolumn{3}{l}{\textbf{Relative to condensed COPOM statement (Treatment 3)}} \\
T4 & -0.646*** & -0.627** \\
& (0.251) & (0.274) \\
\midrule
Demographic Controls & No & Yes \\
Remove Outliers & Yes & Yes \\
\bottomrule
\end{tabular}
\begin{minipage}{\textwidth}
    {\fontsize{8}{8}\selectfont\textit{Notes:} The table reports the probability that subjects' inflation expectation converge to the BCB's target inflation rate of 3% in each treatment group relative to those in the highlighted treatment group. Treatments are described in detail in the text. The second column uses the same specification as the first, but augmented with respondent-specific controls. Results are from Probit regressions. Robust standard errors are reported in parenthesis. * p < 0.1; ** p < 0.05; *** p < 0.01.}
\end{minipage}
\end{table}
